% Source: source_md/intelligence_part2_combined_ja.md
\section{社会的ブロックチェーン — $\Phi_{\text{net}}$ としての分散 $L_E$}


\S{}6〜\S{}8 で、$L_E$ が環境直接相互干渉(アニミズム)、中心化(神)、世代横断的維持(宗教)という形態を取ることを見た。本章は、停止構造を単一権威に依存させずに維持する形態——分散的社会停止——を扱う。これを「社会的ブロックチェーン」と呼ぶ。

まず用語について明確にしておく。本章の「ブロックチェーン」は、暗号通貨や分散台帳の特定の技術を指すものではない。それは、単一権威なしに停止と更新を両立させる構造の特徴を捉えるための、構造的アナロジーである。本論文が借りるのは技術の実装ではなく、その構造的特徴——分散したノードが、中央権威なしに、暫定的合意を形成し、それを更新可能に保つ——だけである。

\bigskip


\subsection{社会的ブロックチェーンの構造的特徴}


社会的ブロックチェーンの中心的特徴は、単一権威なしでの停止と更新の共存である。

\S{}7〜\S{}8 で見た中心化された停止構造(神・宗教)は、停止の曖昧性を低減する代わりに、中心化された参照点の硬直化リスクを負った(\S{}8.5、\S{}8.8)。社会的ブロックチェーンは、このトレードオフに対する別の構造的応答である。それは停止を単一の参照点に集約せず、分散したノード間で維持する。同時に、その停止を固定せず、更新可能に保つ。

停止と更新は、一見、両立しがたい。停止は推論を止めることであり、更新は止めた推論を再び動かすことである。中心化された構造では、この両立は中心点の安定性に依存する。社会的ブロックチェーンは、両立を分散構造そのものによって達成する。個々のノードの停止は暫定的であり、ネットワーク全体としては更新が継続する。

この構造は、\S{}5.8 で論じた「改訂可能性」を、社会的スケールで実装したものである。停止が分散し更新可能であることは、改訂可能性が社会構造に組み込まれていることを意味する。

\bigskip


\subsection{構成要素}


社会的ブロックチェーンの構造を、その構成要素によって記述する。これらは技術的ブロックチェーンの用語からの構造的借用である。

\begin{itemize}
\item \textbf{個体ノード}: 判断を生成する単位。各個体は自らの $L_F$ で推論し、$L_E$ で暫定的停止を生成する。
\item \textbf{ログ}: 永続的な記録。個体ノードの判断とその根拠が記録され、参照可能に保たれる(\S{}5.2 の $L_R$ の社会的拡張)。
\item \textbf{ブロック}: 暫定的に採用された判断。複数のノードによって受容され、当面の停止点として機能する判断の単位。
\item \textbf{フォーク}: 判断の分岐。同一の問題に対して異なる停止が並立する状態。\S{}3.7 の分岐構造の社会的現れ。
\item \textbf{マージ/破棄}: 分岐した判断の統合または放棄。フォークが解消され、停止が再編される過程。
\end{itemize}


これらの要素が、単一権威なしに停止と更新を両立させる機構を構成する。個体ノードが判断を生成し、ログに記録し、ブロックとして暫定採用され、フォークによって分岐し、マージ/破棄によって再編される。この循環が、中央権威なしに継続する。

\bigskip


\subsection{バージョン管理システムとの対比}


社会的ブロックチェーンの性格を明確にするため、技術的なバージョン管理システム(Git のような)と対比する。両者はともに分散・分岐・統合の構造を持つが、その動力学的目的が異なる。

技術的バージョン管理システムは、創造的探索を最大化するよう設計されている。分岐のコストは低く、マージは高速で、ロールバックは容易である。失敗のコストは小さく、多数の実験的分岐を並行させることが推奨される。目的は、可能性の空間を広く探索することである。

社会的ブロックチェーンは、事後安定化を最大化するよう動いている。分岐のコストは高く(社会的分裂を伴う)、マージは遅く(合意形成に時間を要する)、ロールバックは困難である(一度受容された停止の撤回は容易でない)。失敗のコストは大きい。目的は、可能性を広げることではなく、停止を安定化させることである。

この対比は次のように整理される:

\begin{longtable}{@{}p{0.30\linewidth}p{0.30\linewidth}p{0.30\linewidth}@{}}
\toprule
 & 技術的バージョン管理 & 社会的ブロックチェーン \\
\midrule
\endfirsthead
\toprule
 & 技術的バージョン管理 & 社会的ブロックチェーン \\
\midrule
\endhead
動力学的目的 & 創造的探索の最大化 & 事後安定化の最大化 \\
分岐コスト & 低い & 高い \\
マージ速度 & 速い & 遅い \\
ロールバック & 容易 & 困難 \\
失敗の重み & 小さい & 大きい \\
\bottomrule
\end{longtable}


両者のコスト構造は対称的に反転している。これは、技術的システムが探索のために設計されるのに対し、社会的停止構造が安定化のために動いているからである。

\bigskip


\subsection{一貫性ではなく正統性}


技術的ブロックチェーンと社会的ブロックチェーンの最も深い差異は、何を問題とするかにある。

技術的ブロックチェーンは、一貫性 (consistency) の問題を扱う。分散したノードが、矛盾のない単一の状態に合意できるか。これは計算的に定式化可能な問題であり、合意アルゴリズムによって解決される。

社会的ブロックチェーンは、正統性 (legitimacy) の問題を扱う。分散したノードが受容する停止が、正統なものとして成立するか。これは計算的に定式化できない問題である。

正統性は計算できない。それは時間・責任・歴史・分散的受容に依存する。ある停止が正統であるか否かは、アルゴリズムによって決定されるのではなく、ノード間の相互干渉の歴史の中で形成される。正統性は、誰がいつどのように受容したか、その受容がどれだけの時間持続したか、誰が責任を負ったかといった、計算に還元されない要因に依存する。

この差異は本質的である。技術的ブロックチェーンの一貫性問題は、原理的に計算可能である。社会的ブロックチェーンの正統性問題は、原理的に計算不可能である。後者は、$\Phi_{\text{net}}$ における $L_E$ の相互干渉が、計算に還元されない歴史的過程であることを反映している。

この計算不可能性は欠陥ではない。それは社会的停止構造が、固定的アルゴリズムではなく、改訂可能な歴史的過程として停止を維持していることの帰結である。正統性が計算可能であれば、それは固定化され、\S{}8.5 の硬直化リスクに陥る。正統性が計算不可能であることが、社会的停止構造の改訂可能性を保つ。

\bigskip


\subsection{創造は社会で起こらない}


社会的ブロックチェーンについての重要な構造的帰結を述べる。創造は社会では起こらない。創造は個体ノードで起こる。

社会的ブロックチェーンは事後安定化のために動いている(\S{}9.3)。その機能は、生成された判断を受容し、安定化し、改訂可能に保つことである。これは創造の機能ではない。社会は判断を安定化するが、判断を創造しない。

創造——新しい概念、新しい仮説、既存の概念地景にない次元の生成——は、個体ノードで起こる。それは個体の $L_F$ が、しばしば \S{}3.8 の不安定領域すれすれで、既存の準閉包 $\Omega_\theta$ の境界に触れることで生じる。創造は本質的に不安定な過程であり、安定化を目的とする社会構造とは動力学的に相容れない。

ここに、知性の三層動力学における役割分担が現れる。個体ノードは創造する($L_F$ の境界探索)。社会は安定化する($\Phi_{\text{net}}$ の事後受容)。両者は異なる動力学的機能を担い、一方が他方を代替することはできない。社会は創造できず、個体は単独では安定化できない(\S{}3、\S{}5)。

この役割分担は、\S{}11 で論じる不安定領域とも関わる。創造を担う個体ノードは、必然的に不安定領域に接近する。一方、社会はその不安定性を事後的に安定化する。創造と安定化のこの分業が、知性が全体として持続する条件である。

\bigskip


\subsection{三層動力学における位置 — $\Phi_{\text{net}}$ そのものの一形態}


社会的ブロックチェーンを三層動力学の語彙で総括する。

社会的ブロックチェーンは、$\Phi_{\text{net}}$ そのものの一形態である。より正確には、$\Phi_{\text{net}}$——個体間 $L_E$ の相互干渉ネットワーク——を、明示的な分散構造として運用する形態である。

\S{}6〜\S{}8 の形態と対比すると、その位置が明確になる。アニミズムは $L_E$ を環境に分散させた。神は $L_E$ を中心化した。宗教は中心化された $L_E$ を世代横断的に維持した。社会的ブロックチェーンは、$L_E$ ネットワークそのものを、中心化せずに、明示的な分散構造として運用する。

\begin{figure}[H]
\centering
\includegraphics[width=\linewidth]{fig7_evolution_of_stopping_structures_ja.pdf}
\caption{停止構造の進化。アニミズム・神・宗教・科学・社会的ブロックチェーンは、同一の創発停止層 $L_E$ を安定化・移動・共有するための異なる操作である。}
\label{fig:part2-stopping-evolution}
\figurenote{五つの形態は、設置・圧縮・維持・改訂・同期という停止操作として配置される。これは進歩物語ではなく、$L_E$ の操作レパートリーの多様化である。}{停止は自由の否定ではない。歴史的停止構造は、開かれた推論を有限な系の中で持続可能にするための運用形式として読める。}
\end{figure}

この系列を停止構造の運用史として読めば、さらに位置づけは明確になる。宗教が停止構造の長期維持を制度化した運用形態であるならば、科学は停止構造の更新を制度化した運用形態として理解できる。社会的ブロックチェーンはさらにその上で、停止構造の分散的共有と同期を担う仕組みとして位置づけられる。

重要なのは、これがアニミズムの分散への「回帰」ではないことである(\S{}8.8 で予告)。アニミズムの分散は、個体と環境の直接相互干渉であった。社会的ブロックチェーンの分散は、個体間 $L_E$ の相互干渉ネットワークの運用である。前者は個体−環境の関係、後者は個体−個体の関係である。両者はともに分散的だが、分散している対象が異なる。

この構造は、Bio-IFGT における「個体間相互拘束による準閉包維持」と直接相同である。Bio-IFGT は、生命の準閉包が個体間の相互拘束によって維持されることを示した。社会的ブロックチェーンは、知性の停止構造($L_E$)が個体間の相互干渉($\Phi_{\text{net}}$)によって維持されることを示す。生命層の個体間準閉包維持と、知性層の個体間停止維持は、同一の構造の異なる層における現れである。

すなわち、個々の $L_E$ は個体内に閉じない。それはノード間相互干渉によって相互に維持される。個体の $L_R$(ログ)が社会的に共有・参照されることで、$L_E$ の分散維持が可能になる。社会的ブロックチェーンは、この分散維持の明示的な構造である。

\bigskip


\subsection{次章への橋渡し}


\S{}6〜\S{}9 で、$L_E$ と $\Phi_{\text{net}}$ が取る歴史的・社会的形態——アニミズム・神・宗教・社会的ブロックチェーン——を分析した。これらに共通するのは、停止が正しさによってではなく、構造的安定性と改訂可能性によって維持されるという点である。

社会的ブロックチェーンにおいて、これは特に明確に現れる。正統性は計算できず(\S{}9.4)、停止は真理によってではなく分散的受容によって成立する。ここから、停止と真理の関係についての根本的な問いが生じる。社会的停止が真理を保証しないなら、知識とは何か。正しく信じられ正当化された判断が、なお知識でないことがあるとすれば、知識を成立させているのは何か。

次章 \S{}10 で扱うゲティア問題は、この問いに対する構造的診断である。ゲティア問題は、$C$(閉包度)が真理測度ではなく安定性測度であること——$L_E$ が正しい停止ではなく改訂可能な停止を供給すること——の、認識論における証拠である。

\bigskip


\subsection{要約}


本章で示したのは次の点である:

\begin{enumerate}
\item 「社会的ブロックチェーン」は暗号通貨技術ではなく、単一権威なしに停止と更新を両立させる構造のアナロジーである
\item その構成要素は個体ノード・ログ・ブロック・フォーク・マージ/破棄である
\item 技術的バージョン管理(創造的探索の最大化)と社会的ブロックチェーン(事後安定化の最大化)はコスト構造が対称的に反転している
\item 技術的ブロックチェーンは一貫性問題を、社会的ブロックチェーンは正統性問題を扱う。正統性は計算できない
\item 正統性の計算不可能性は欠陥ではなく、改訂可能性を保つ条件である
\item 創造は社会では起こらず個体ノードで起こる。社会は安定化を担う。創造と安定化は役割分担される
\item 社会的ブロックチェーンは $\Phi_{\text{net}}$ を明示的分散構造として運用する形態であり、Bio-IFGT の個体間準閉包維持と相同である
\item 停止が真理でなく構造的安定性によることが、\S{}10 のゲティア問題への問いを生む
\end{enumerate}


次章 \S{}10 では、$C$ が真理測度ではなく安定性測度であることを、ゲティア問題を通じて構造的に診断する。

\bigskip
